\clearpage
\pagestyle{plain}
\singlespacing

{\crsTOCheading СОДЕРЖАНИЕ\par\vspace{1ex}}

\crsTOCline[0pt]{\normalfont\textbf{РЕФЕРАТ}}{2}
\crsTOCline[0pt]{\normalfont\textbf{ВВЕДЕНИЕ}}{4}
\crsTOCline[0pt]{\normalfont\textbf{ГЛАВА\,1}. Теоретические основы генерации случайных чисел}{6}
\crsTOCline[1.5em]{\normalfont{}1.1. Основные понятия и определения}{6}
\crsTOCline[1.5em]{\normalfont{}1.2. Классификация генераторов случайных чисел}{8}
\crsTOCline[1.5em]{\normalfont{}1.3. Методы оценки случайности последовательностей}{10}
\crsTOCline[0pt]{\normalfont\textbf{ГЛАВА\,2}. Анализ физических источников энтропии и их применение в генераторах истинно случайных последовательностей}{15}
\crsTOCline[1.5em]{\normalfont{}2.1. Электронные источники шума}{15}
\crsTOCline[1.5em]{\normalfont{}2.2. Оптические и квантовые источники}{24}
\crsTOCline[1.5em]{\normalfont{}2.3. Механические и акустические источники}{28}
\crsTOCline[1.5em]{\normalfont{}2.4. Постобработка}{32}
\crsTOCline[1.5em]{\normalfont{}2.5. Уязвимости и атаки}{35}
\crsTOCline[0pt]{\normalfont\textbf{ГЛАВА\,3}. Практический анализ и сравнение источников энтропии}{39}
\crsTOCline[1.5em]{\normalfont{}3.1. Сравнительный анализ источников энтропии}{39}
\crsTOCline[1.5em]{\normalfont{}3.2. Архитектура экспериментального стенда и программная инфраструктура}{41}
\crsTOCline[1.5em]{\normalfont{}3.3. Аналоговый усилительный тракт}{47}
\crsTOCline[1.5em]{\normalfont{}3.4. Тепловой шум резистора}{55}
\crsTOCline[1.5em]{\normalfont{}3.5. Лавинный шум обратносмещённого BE-перехода биполярного транзистора}{60}
\crsTOCline[1.5em]{\normalfont{}3.6. Акустический шум электретного микрофона}{75}
\crsTOCline[1.5em]{\normalfont{}3.7. Электромагнитные наводки на плавающий вход АЦП}{77}
\crsTOCline[1.5em]{\normalfont{}3.8. MEMS-инерциальный датчик MPU-6050}{78}
\crsTOCline[1.5em]{\normalfont{}3.9. Дрожание тактовых сигналов (clock jitter)}{79}
\crsTOCline[1.5em]{\normalfont{}3.10. Итоги практической части и сравнение источников}{80}
\crsTOCline[1.5em]{\normalfont{}3.11. Перспективные направления развития физических генераторов случайных последовательностей (ФГСП)}{83}
\crsTOCline[0pt]{\normalfont\textbf{ЗАКЛЮЧЕНИЕ}}{85}
\crsTOCline[0pt]{\normalfont\textbf{СПИСОК ИСПОЛЬЗОВАННЫХ ИСТОЧНИКОВ}}{87}
\crsTOCline[0pt]{\normalfont\textbf{ПРИЛОЖЕНИЯ}}{89}
\crsTOCline[1.5em]{\normalfont Приложение~1\,--- Статистические тесты пакета NIST STS}{89}
\crsTOCline[1.5em]{\normalfont Приложение~2\,--- Сравнительная таблица характеристик источников энтропии}{93}
